% =========================================================
% FRAGMENTO 3 — AFLORAMENTOS, GRANULOMETRIA, TEMPO
% GEOLÓGICO, LYELL, BOWEN, CORES DOS MINERAIS E WENTWORTH
%
% Este arquivo contém somente o corpo dos verbetes.
% Usa os comandos já definidos no arquivo principal:
% \letra, verbete, \termo, \relacao, \vertambem e pensador.
%
% OBSERVAÇÃO: Georges Cuvier já aparece no fragmento anterior.
% Não foi repetido aqui para evitar dois verbetes iguais.
%
% Grafia correta: Siccar Point.
% =========================================================





% =========================================================
% LETRA M — inserir depois de "Minerais e sólidos amorfos".
% NÃO repetir \letra{M}.
% =========================================================



% =========================================================
% LETRA R — inserir dentro da letra R, de preferência antes
% de "Rochas sedimentares: tipos principais".
% NÃO repetir \letra{R}.
% =========================================================



\begin{verbete}{Siccar Point}
\termo{Siccar Point} é um afloramento costeiro da Escócia famoso por apresentar uma discordância angular reconhecida por James Hutton em 1788. No local, camadas quase verticais de arenitos grauváquicos, siltitos e lamitos do Siluriano são recobertas de forma discordante por brechas e arenitos avermelhados mais jovens do \textit{Old Red Sandstone}, depositados com inclinação muito menor.

A geometria observada exige uma longa sucessão de acontecimentos:

\begin{enumerate}
  \item deposição e litificação dos sedimentos silurianos;
  \item deformação, dobramento e basculamento das camadas antigas;
  \item soerguimento e erosão, que truncaram essas camadas;
  \item nova deposição de sedimentos continentais sobre a superfície erosiva;
  \item novo soerguimento e erosão, expondo a discordância atual.
\end{enumerate}

\begin{center}
\begin{tikzpicture}[x=0.72cm,y=0.72cm,every node/.style={font=\sffamily\scriptsize,align=center}]
  \path[use as bounding box] (-0.2,-0.45) rectangle (8.65,4.85);

  % Camadas antigas, quase verticais
  \foreach \x in {0.15,0.75,1.35,1.95,2.55,3.15,3.75,4.35,4.95,5.55,6.15,6.75,7.35,7.95} {
    \draw[uerjblue,thick] (\x,0.15) -- ({\x+0.75},2.35);
  }
  \fill[white] (-0.2,2.30) -- (8.65,2.30) -- (8.65,4.85) -- (-0.2,4.85) -- cycle;
  \draw[uerjred,very thick] (-0.05,2.30) -- (8.45,2.30);

  % Camadas jovens, pouco inclinadas
  \foreach \y in {2.65,3.10,3.55,4.00} {
    \draw[uerjgold,thick] (0.05,{\y+0.30}) -- (8.45,\y);
  }

  \node[fill=white,inner sep=1.5pt] at (2.10,1.25) {camadas silurianas\\quase verticais};
  \node[fill=white,inner sep=1.5pt] at (6.25,3.55) {camadas mais jovens\\pouco inclinadas};
  \node[fill=white,inner sep=1.5pt,text=uerjred] at (4.25,2.30) {superfície de discordância};
\end{tikzpicture}
\end{center}

Siccar Point forneceu uma demonstração visual do \termo{tempo profundo}: cada etapa exigiu tempo para ocorrer, e a superfície de discordância representa um intervalo não preservado diretamente como camada sedimentar.

\relacao{Siccar Point relaciona afloramento, discordância angular, ciclos de deposição e erosão e a interpretação do tempo geológico profundo.}

\vertambem{Hutton, James, discordância angular, afloramento, registro estratigráfico, tempo geológico}
\end{verbete}


